% WORKING GUIDE: extended after each completed part of the cote24 audit.
\subsection*{Piotr Blass's translation: the opening and its notation}
The English title page identifies the work as ``Stage 0 -- translation by Piotr Blass,'' and the following Summary places the translated \emph{Commentaires} before the contents. Its queries belong to this witness: ``reorganize???'' corresponds to \emph{réorganisation de fond en comble} on French site 2; the typed question about no.\,16 and its handwritten ``Ask AG.'' refer to the French statements on sites 2--3. In Section 1, the English develops the same quotient-Module, hyperplane, relative-divisor, and incidence constructions as French sites 4--8. The opaque strips on English sites 4 and 6 remain unreadable, with their positions identified rather than completed from the French. English site 5 asks which paragraph 17 is meant and queries the separate circle above the functor domain; both corresponding passages are on French site 5. English site 7 retains four unfilled arrow positions and asks whether \emph{Module} always means a coherent or quasi-coherent sheaf of modules, at the passage corresponding to French site 7. Its drawn diagram contains four connectors without arrowheads. The comparison also records exact differences of notation: the final term of the transpose formula on English site 4 is $\underline O_P(1)$, while French site 4 types $\underline D_P(1)$ and then names a section of $\underline L_P(1)$; English site 4 types an ordinary $p*$ where the French has a handwritten subscript star; English site 5 lacks the French dual marks in the symmetric-power expression and on the final projective space; English site 6 lacks one dual mark in the section over $S'$; and English site 7 types the continuous string $XPV$ where French site 7 says that $F$ is a Module on $X$. These visible forms are preserved. Full-height baseline letters and genuinely raised characters have been compared separately, since the typescript uses both. Sources: \href{https://grothendieck.umontpellier.fr/web/viewer.html?file=../24.pdf\#page=2}{cote 24, sites 1--8}; the corresponding French passages are \href{https://grothendieck.umontpellier.fr/web/viewer.html?file=../25.pdf\#page=3}{cote 25, sites 2--8}.

\paragraph{Blass's Section 2 (sites 9--19 opening).}
The translation follows the local study from irreducibility and finiteness through associated cycles, regularity, smoothness, and separability, and retains the author's comments on reorganizing the results. Its queries can be located precisely: site 9 questions the insertion \emph{En tous cas} in French 8; site 10 asks about underlined capitals; site 12 queries the added alternative in Proposition 2.6 and the paragraph-3 reference of French 9--10; site 13 leaves the marginal remark about the \emph{dim.\,|\,dém.} of 2.1 unread and retains four reference placeholders from French 11; site 15 leaves the argument of $k(\ )$ blank, as French 13 does; and site 17 asks about the two paragraph-6 references on French 13--14. Both opaque masks on English 16 remain unreadable. The first ``PB incomprehensible'' on English 18 corresponds to the repeated $Z_i$ in French 14's first assertion of 2.16; the repetition is present in both witnesses. Its second query accompanies the reduction to an algebraically closed field, whose order relative to the separability sentence differs from the French insertion. The mathematical comparisons retain all observed differences: English 9 inserts $k$ in $\operatorname{Spec}k(\eta)$, where French 8 has $\operatorname{Spec}(\eta)$; English 10 types $Z_\eta$ in the proof where French 8 corrects the letter to $Y$; English 11 types $Z1$ where French 9 has $Z'$; English 13 types a capital $N$ in the coprofundity bound and $H_\eta$ in the conclusion, where French 11 has lowercase $n$ and the doubtful $E|F$ with indices $\eta,y$. English 14 places its first bar over \emph{dim}, omits the inverse-image exponent in $h(Z)$, and says that constructibility is already known by applying 2.3, while French 11--12 bars $f(Z)$, writes the inverse-image exponent, and says that constructibility must be known to apply 2.3. English 15 has a lowercase $y$ and ordinary arrow into $U_{k(\eta)}$ against French 12's capital $Y$ and hooked arrow, and retains the baseline $1$ and lowercase $u/k$ in its differential notation. English 16 places the bound $1\leq i\leq n$ beneath the quotient sum; French 13 places it after the inserted generators $a_i$. English 16--17 omit the French marginal qualification beside 2.14, whose first term remains doubtful, while English 17 types $k(y)$ in the equivalence and omits ``geometrically'' before the first ``normal'' in 2.15. Capital $K(\eta)$ occurs on English 16--17 where the corresponding French uses lowercase $k$. Several additional underlines on $Y$ are visible on English 9, 11, 12, and 14. English 18 replaces the French ranges 2.13--2.17 and 2.12--2.17 by conjunctions of their endpoints; the explanation continues on English 19, with an extra typed arrow before ``smooth = regular.'' None of these forms is repaired from the other witness. Sources: \href{https://grothendieck.umontpellier.fr/web/viewer.html?file=../24.pdf\#page=10}{cote 24, sites 9--19}; \href{https://grothendieck.umontpellier.fr/web/viewer.html?file=../25.pdf\#page=9}{cote 25, sites 8--15}.

\paragraph{Blass's Section 3 (sites 19--21).}
The translation presents the Bertini--Zariski theorem through the function-field description and Zariski's lemma, retains the request that Dieudonné choose a proof from the literature, and follows with the geometric-integrality and connectedness corollaries. The numbering query on site 19 corresponds to the overwritten heading on French 15; ``[PB: check this!]'' points to the inclusion described by sending $S_i$ to $S_i$ on French 16. Several literal formulas require care. English 19 has $\dim\overline{f(X)}\geq0$ where French 16 has $>0$, and ends $L=K(S_1,\ldots,S_n)$ at $n$ where the French ends at $n-1$. English 20 uses $a_iS_1$ where French 16 has $a_1S_1$, and omits the minus before $t_it_n^{-1}$ in the definition of $a_i$. Its lemma omits the French definition $L=K(S_1,\ldots,S_m)$, while the sum has printed capital limits $I$ and $M$ against French 17's $1$ and $m$; the English field name is the literal baseline $k1$. English 21 uses a clear typed $\sigma$ and $f(X)$ where the checked French 17 retains doubts about the script family letter and the argument $Z|X$. It also prints capital $K(\eta)$ in the conclusion and ``$X_0$ and $X'$'' at the start of the chain, against the French lowercase field letter and equality. Its handwritten query points to the paragraph-6 reference accompanying Hartshorne's theorem. The guide reports these readings and retains the French uncertainties. Sources: \href{https://grothendieck.umontpellier.fr/web/viewer.html?file=../24.pdf\#page=20}{cote 24, sites 19--21}; \href{https://grothendieck.umontpellier.fr/web/viewer.html?file=../25.pdf\#page=16}{cote 25, sites 15--18}.
